\documentclass[runningheads]{llncs}

\usepackage{graphicx}
\usepackage{amsmath}
\usepackage{amssymb}
\usepackage{booktabs}
\usepackage{xcolor}
\usepackage{hyperref}
\usepackage{cite}
\usepackage{orcidlink}
\renewcommand\UrlFont{\color{blue}\rmfamily}

\begin{document}

\title{Adaptive Preprocessing Policy Search via Reinforcement Learning
for Noisy Multiclass Spectral Classification}

\titlerunning{Adaptive Preprocessing Policy Search via RL}

\author{Natalya Smith\orcidlink{0000-0000-0000-0000}}
\authorrunning{N. Smith}
\institute{Independent Researcher, Auckland, New Zealand\\
\email{mnb8479@autuni.ac.nz}}

\maketitle

%-----------------------------------------------------------------------
\begin{abstract}
% PLACEHOLDER — complete after experiments.
%
% Target structure (≤200 words):
%   1. Problem: fixed preprocessing pipelines for Raman spectra cannot
%      adapt to the learning dynamics of the classifier or the noise
%      structure of the measurement.
%   2. Gap: existing work on adaptive preprocessing uses heuristic
%      selection; the validity of the evaluation metric used to guide
%      selection has not been examined.
%   3. Approach: RL agent selects among N discrete preprocessing
%      pipelines per fold, rewarded by the chosen evaluation metric.
%      Three reward variants compared: macro-F1, balanced accuracy,
%      ROC-AUC (one-vs-rest).
%   4. Key finding: [INSERT AFTER EXPERIMENTS]
%   5. Key claim: preprocessing selection should be treated as a
%      measurement validity problem, not a hyperparameter search.

\textbf{[Placeholder — complete after experimental results.]}

\keywords{Raman spectroscopy \and Drug verification \and Reinforcement
learning \and Preprocessing policy search \and Evaluation validity \and
Multiclass classification}
\end{abstract}

%-----------------------------------------------------------------------
\section{Introduction}

% WRITING NOTES:
% Open with the measurement problem, not the ML problem.
% Raman spectroscopy produces overlapping spectral fingerprints.
% The preprocessing pipeline determines what features the classifier sees.
% A pipeline chosen once, before training, cannot adapt to the noise
% structure of individual measurement sessions.
%
% Then: the RL reward signal is a measurement instrument — the same
% validity argument as Paper 1, but now for a harder problem:
%   - multiclass (not binary)
%   - the policy object is preprocessing, not augmentation
%   - the cost of confusion is asymmetric (drug verification context)
%
% Research questions (two, explicit):
%   Q1: Can RL-guided preprocessing policy search outperform a fixed
%       pipeline on multiclass Raman spectral classification?
%   Q2: Does the choice of reward metric (macro-F1, balanced accuracy,
%       ROC-AUC) affect the quality of the learned preprocessing policy?
%
% Contributions (as findings, not activities):
%   1. Framework transfer: the adaptive policy search framework from
%      Paper 1 (image augmentation) is re-instantiated with a different
%      policy object (spectral preprocessing) and a harder evaluation
%      problem (multiclass, asymmetric costs).
%   2. Reward metric comparison for spectral preprocessing.
%   3. Empirical demonstration that preprocessing choice is consequential
%      for multiclass spectral classification under measurement noise.

\textbf{[Introduction — draft after experimental design is finalised.]}

The primary contributions of this work are:
\begin{itemize}
  \item \textbf{[CONTRIBUTION 1]} % preprocessing policy search framework
  \item \textbf{[CONTRIBUTION 2]} % reward metric comparison
  \item \textbf{[CONTRIBUTION 3]} % instrument validity argument for spectra
\end{itemize}

%-----------------------------------------------------------------------
\section{Related Work}

\subsection{Raman Spectroscopy and Preprocessing}

% Cover:
% - Why Raman preprocessing matters (overlapping fingerprints, noise)
% - Standard pipelines: SNV, derivatives, baseline correction
% - Evidence that pipeline choice affects classification performance
% - Gap: adaptive vs static preprocessing

\textbf{[Related work — draft after literature review.]}

\subsection{Adaptive Policy Search for Data Processing}

% Cover:
% - AutoAugment (Cubuk 2019) and PBA (Ho 2019) — image augmentation
% - Paper 1 (this work's predecessor) — augmentation in histopathology
% - Any existing work on adaptive spectral preprocessing (likely sparse)
% - The gap: reward metric validity in policy search has not been examined

\subsection{Evaluation Metric Validity in Applied ML}

% This subsection carries over directly from Paper 1 and extends it.
% Key addition: in multiclass settings, metric choice is more consequential
% because there is no natural symmetric loss — different compounds have
% different misclassification costs in a drug verification context.
% Macro-F1, balanced accuracy, and OvR-AUC encode different assumptions
% about class importance. This is a construct validity question.

%-----------------------------------------------------------------------
\section{Methodology}

\subsection{Dataset and Task}

% PLACEHOLDER — fill in after you confirm what data you are using.
%
% Describe:
% - Source of Raman spectra (spectral-drug-verification repo data)
% - Number of spectra, number of compounds (classes)
% - Wavenumber range retained (150–3425 cm^{-1})
% - Class distribution and any imbalance
% - Important: state that this is a prototype dataset, not production
%   analyser data (honest about external validity limits)
%
% Note the key methodological difference from BreakHis:
% No patient-level leakage concern — each spectrum is independent.
% Standard stratified k-fold is appropriate here.

\textbf{[Dataset description — confirm data and fill in.]}

\subsection{Preprocessing Policy Action Space}

% Describe the 9 discrete preprocessing pipelines (A0–A8).
% Table here:

\begin{table}[t]
\caption{Preprocessing pipeline action space.}\label{tab:actions}
\centering
\begin{tabular}{lll}
\toprule
Action & Pipeline & Description \\
\midrule
$A_0$ & Raw & No preprocessing \\
$A_1$ & SNV & Standard Normal Variate \\
$A_2$ & Baseline & Polynomial baseline correction (degree 3) \\
$A_3$ & SG & Savitzky-Golay smoothing (window 11, poly 3) \\
$A_4$ & 1st deriv. & First derivative via SG \\
$A_5$ & 2nd deriv. & Second derivative via SG \\
$A_6$ & SNV + 2nd & SNV followed by second derivative [\emph{static baseline}] \\
$A_7$ & SNV + SG & SNV followed by SG smoothing \\
$A_8$ & Window & Restricted wavenumber window (500--2000 cm$^{-1}$) \\
\bottomrule
\end{tabular}
\end{table}

% Note the important difference from Paper 1:
% In Paper 1, augmentation is applied stochastically per sample during
% training. Here, the preprocessing pipeline is applied deterministically
% to all spectra in the fold (train, val, and test). This means:
% (a) the policy is more consequential — it determines what the
%     classifier sees at every step, including evaluation.
% (b) there is no stochasticity to average out — a bad pipeline
%     consistently harms the classifier.
% (c) the pipeline applied during search must also be applied at test
%     time, which is enforced in the implementation.

\subsection{RL Framework}

% MDP formulation — same structure as Paper 1 but with different objects.
%
% State: (last_action_id, metric_bucket)
%   metric_bucket: 0=low(<0.85), 1=mid(0.85-0.95), 2=high(>0.95)
%
% Action: preprocessing pipeline id in {0,...,8}
%
% Reward: delta_metric after retraining classifier under new pipeline
%
% Three reward variants:
%   macro-F1: recommended — class-averaged, penalises poor performance
%             on any single compound equally regardless of class size
%   balanced accuracy: average recall per class — simpler but less
%             sensitive to precision differences
%   OvR-AUC: threshold-independent, captures discrimination globally
%
% Explain why macro-F1 is argued to be most valid for this context:
%   In drug verification, identifying the wrong compound is dangerous
%   regardless of which compound it is. A metric that averages over
%   classes with equal weight encodes this concern; accuracy (which
%   weights by class size) does not.

The augmentation selection problem is formalised as an MDP
$\mathcal{M} = (S, A, P, R, \gamma)$:

\begin{itemize}
  \item \textbf{State} $S$: $(a_t, b_t)$ where $a_t$ is the last
        preprocessing action and $b_t \in \{0,1,2\}$ is the
        discretised validation metric bucket.
  \item \textbf{Action} $A$: preprocessing pipeline $a \in \{0,\ldots,8\}$
        (Table~\ref{tab:actions}).
  \item \textbf{Reward} $R$: $r_t = \text{metric}_t - \text{metric}_{t-1}$,
        where metric is one of macro-F1, balanced accuracy, or OvR-AUC.
  \item \textbf{Discount} $\gamma = 0.9$.
\end{itemize}

The agent uses $\epsilon$-greedy tabular Q-learning with the update:
\begin{equation}
  Q(s,a) \leftarrow Q(s,a) + \alpha
  \bigl[r + \gamma \max_{a'} Q(s',a') - Q(s,a)\bigr]
\end{equation}
with $\alpha=0.1$, $\epsilon$ decaying from 0.3 to 0.05.

\subsection{Reward Metric Validity Argument}

% This is the core methodological section — do not bury it.
% Argument structure:
%   1. In multiclass spectral classification, metric choice encodes
%      assumptions about class importance and error costs.
%   2. Macro-F1 treats all classes equally — appropriate when all
%      compound misidentifications carry similar risk.
%   3. Balanced accuracy is simpler but ignores precision.
%   4. OvR-AUC is threshold-independent but requires the model to
%      output calibrated probabilities.
%   5. The choice of reward metric is therefore a construct validity
%      decision: which metric best operationalises the target construct
%      (correct compound identification under measurement noise)?
%   6. We test this empirically rather than assuming.

\textbf{[Reward validity argument — draft after confirming class structure.]}

\subsection{Classifier}

% Random Forest as default for spectral data.
% Justify: RF handles high-dimensional, correlated features well;
% does not require normalisation; provides feature importance estimates
% that can be used to validate whether the selected preprocessing
% pipeline exposes informative wavenumber regions.
%
% Note: the framework is classifier-agnostic. Results with SVM or
% a 1D CNN can be reported as secondary comparisons.

\subsection{Experimental Conditions}

\begin{enumerate}
  \item \textbf{No preprocessing} ($A_0$): raw spectra.
  \item \textbf{Static baseline} ($A_6$): SNV + second derivative
        (current pipeline in \texttt{spectral-drug-verification}).
  \item \textbf{RL-macro-F1}: adaptive policy, reward = $\Delta$macro-F1.
  \item \textbf{RL-balanced-acc}: adaptive policy, reward = $\Delta$balanced accuracy.
  \item \textbf{RL-OvR-AUC}: adaptive policy, reward = $\Delta$OvR-AUC.
\end{enumerate}

\subsection{Evaluation}

% Metrics: macro-F1, balanced accuracy, per-class recall, confusion matrix,
% OvR-AUC, log loss (calibration proxy).
%
% Emphasise per-class recall — in drug verification, knowing which
% specific compound is misidentified matters. A confusion matrix that
% clusters errors between chemically similar compounds is a different
% failure mode than one that scatters errors randomly.

%-----------------------------------------------------------------------
\section{Results and Discussion}

% PLACEHOLDER — complete after experiments.
%
% Structure:
%   4.1 Effect of preprocessing condition
%       Table: no preprocessing / static / RL-macro-F1
%       Claim: adaptive pipeline outperforms static?
%
%   4.2 Reward metric comparison
%       Table: RL-macro-F1 / RL-balanced-acc / RL-OvR-AUC
%       Key result: which metric produces the best recall on hard classes?
%       Is the pattern consistent with Paper 1's finding?
%
%   4.3 Per-class analysis
%       Confusion matrices — which compounds are hard to separate?
%       Does the chosen pipeline help with the hardest class pairs?
%
%   4.4 Connection to Paper 1
%       What transfers? What changes?
%       The reward-metric finding: does macro-F1 dominate in the
%       multiclass case the way AUC dominated in the binary case?

\textbf{[Results — complete after experimental runs.]}

%-----------------------------------------------------------------------
\section{Conclusion}

% Two contributions to state:
%   1. The adaptive policy search framework transfers from image
%      augmentation (Paper 1) to spectral preprocessing, demonstrating
%      that the policy-learning abstraction generalises across
%      measurement modalities.
%   2. The reward metric comparison confirms / extends the finding
%      from Paper 1: the choice of evaluation instrument used to guide
%      policy search is consequential, and this is a construct validity
%      question that warrants empirical investigation.
%
% Limitations:
%   - Single prototype dataset (not production analyser data)
%   - Random Forest as default classifier (not necessarily optimal)
%   - Policy is per-fold, not adaptive within a fold
%
% Future work:
%   - Production Raman data
%   - Concentration estimation as a second objective
%   - Abstention / uncertainty-aware reward

\textbf{[Conclusion — complete after results.]}

%-----------------------------------------------------------------------
\bibliographystyle{splncs04}
\begin{thebibliography}{20}

% ANCHOR CITATIONS (do not remove)

\bibitem{cubuk2019}
Cubuk, E.D., Zoph, B., Mane, D., Vasudevan, V., Le, Q.V.: AutoAugment:
learning augmentation strategies from data. In: IEEE CVPR, pp.~113--123 (2019)

\bibitem{ho2019}
Ho, D., Liang, E., Chen, X., Stoica, I., Abbeel, P.: Population based
augmentation: efficient learning of augmentation policy schedules. In:
ICML, vol.~97, pp.~2731--2741. PMLR (2019)

\bibitem{hospedales2021}
Hospedales, T., Antoniou, A., Micaelli, P., Storkey, A.: Meta-learning
in neural networks: a survey. IEEE TPAMI \textbf{44}(9), 5149--5169 (2021)

\bibitem{maleki2022}
Maleki, F., Ovens, K., Gupta, R., Reinhold, C., Spatz, A., Forghani,
R.: Generalisability of machine learning models: quantitative evaluation
of three methodological pitfalls. Radiology: AI \textbf{5}(1), e220028 (2022)

% ADD: Raman preprocessing references
% ADD: Drug verification / spectroscopy ML references
% ADD: Multiclass evaluation metric references
% ADD: Paper 1 self-citation once on arXiv

\end{thebibliography}

\end{document}
